\documentclass[12pt]{article}
\usepackage[utf8]{inputenc}
\usepackage[T1]{fontenc}
\usepackage{hyperref}
\title{Tawara Language Resources for K-12 Curriculum Development}
\author{Compiled by Claude Sonnet 5 (claude-sonnet-5) for Andr\'as Kornai}
\date{}
\begin{document}
\maketitle

\section{Introduction}
Tawara (ISO 639-3 twl; also Tavara, Tawala) is spoken in Mozambique and neighbouring areas, with dialects Tawara-Chioco and Tawara-Daque reported by Joshua Project. Population figures conflict across sources (50,000 in one, 178,000 speakers in another). An older source treats Tawara as a variant of Shona-Korekore. Very little was found, and nothing for schools.

\section{Primary Educational Resources}
None found.

\section{Linguistic and Reference Materials}
A Leiden University publication lists Tawara as a dialect variant of Shona-Korekore, which differs from treating it as a separate language: \url{https://scholarlypublications.universiteitleiden.nl/access/item%3A2938467/view}

\section{Religious and Community Resources}
Joshua Project lists a New Testament (1985) and no complete Bible: \url{https://joshuaproject.net/languages/twl}. Its Tawara in Mozambique page is at \url{https://joshuaproject.net/index.php/people_groups/15302/MZ}.

\section{Digital and Audio Resources}
Global Recordings Network has a page for twl, which suggests audio recordings exist: \url{https://globalrecordings.net/en/language/twl}

\section{Curriculum Development Context}
Documentation is thin and its classification is unsettled. A curriculum would need fieldwork and a decision on dialect and orthography. Portuguese is the usual language of instruction in Mozambique.

\end{document}